\section{Introduction}
Emerging studies on the generalizability of deep neural networks have revealed that the existing models, which assume independent and identically distributed (i.i.d.) training and test distribution, are not robust to significant distribution shifts between training and test distribution, \eg, backgrounds \cite{xiao2020background_challenge_bgc}, geographic distribution \cite{de2019does}, demographic statistics \cite{yang2020towards, scimeca2022wcst-ml}, textures \cite{geirhos2019cnn_biased_towards_texture, bahng2019rebias}, or day-to-night shifts \cite{dai2018dark, michaelis2019benchmarking}.
Domain generalization (DG) aims to learn robust representations against distribution shifts from multiple source {domains} during training. The trained model is evaluated on an unseen domain to measure the robustness.
% Domain generalization (DG) problem measures the robustness against the distribution shifts by evaluating the model on an unseen domain during training.
The existing DG approaches have tried to learn invariant features across multiple {domains} \cite{ganin2016dann,arjovsky2019irm,li2018mldg,cha2021swad,sun2016coral,bui2021mdsdi,zhou2021mixstyle}.
%by minimizing feature divergences between the source domains~\cite{muandet2013icml_DIFL,ganin2016dann,li2018mmd,sun2016coral,matsuura2020aaai_mixture_mld,li2018cdann,zhao2020er_entropy_regularization}, simulating domain shifts based on meta-learning~\cite{li2018mldg,balaji2018metareg,dou2019masf,li2019episodic_epi-fcr,zhang2020arm}, robust optimization \cite{arjovsky2019irm, krueger2020vrex, Sagawa2020GroupDRO, shi2022gradient}, or augmenting source domain examples~\cite{shankar2018crossgrad,nuriel2020padain,zhou2021mixstyle,nam2019sagnet,carlucci2019jigsaw_jigen,zhou2020l2a_ot,bai2020decaug}.
However, recent studies \cite{gulrajani2020domainbed, koh2021wilds} have shown that simple baselines without learning invariant features are comparable to or even outperform the existing DG methods on the diverse DG benchmarks with a fair evaluation protocol in realistic settings (\eg, using ResNet-50 instead of ResNet-18 \cite{he2016_cvpr_resnet}).
%(\eg, from ResNet-18 to ResNet-50 \cite{he2016_cvpr_resnet}). 
We presume that it is because training and test distributions differ too significantly to learn domain-invariant features by the training distribution only.

Instead of learning domain-invariant features, we let a model learn similar features to ``oracle'' representations, \ie, an optimal model generalized to \emph{any} domain. In particular, we re-formulate the DG problem by maximizing the mutual information (MI) between the oracle model representations and the target model representations while preserving the training loss on source domains. However, the oracle model is not achievable in practice. Hence, we use a large pre-trained model (\eg, ImageNet pre-trained ResNet-50 \cite{he2016_cvpr_resnet}) as an approximation. With this approximation, we derive a tractable variational lower bound of the proposed maximization problem, named \fullname{} (\method{}). At a high level, our \method{} objective consists of two objectives: an original target task (\ie, an ERM objective) and a regularization term between the pre-trained model and the current target model. Note that the standard DomainBed benchmark \cite{gulrajani2020domainbed} uses the ImageNet pre-trained ResNet-50 as the initialization of a DG method, thus, we use the pre-trained ResNet as the initialization and the approximation of the oracle model at the same time.

While a naive fine-tuning approach of a large pre-trained model can harm the robustness against distribution shifts \cite{wortsman2021robust, kumar2022iclr_finetuning_distort}, our proposed algorithm remarkably improves the robustness against unseen domains during fine-tuning in a plug-and-play manner to any scale of the backbone model and datasets. 
%In our experiment, we observe that the naive fine-tuning of large pre-trained models can perform worse than the ERM baseline with ImageNet pre-trained model. For example, ERM with ImageNet the pre-trained ResNet shows 64.2\% average accuracy, while ERM with the pre-trained CLIP ViT shows 61.1\%, an inferior accuracy than ImageNet pre-trained model. 
% 1. large-scale pretrained model 이 더 낮은 성능을 내는 경우가 있음. e.g. IN-ERM vs. CLIP-ViT
% 2. 우리 모델은 다양한 스케일에서 성능을 향상시킴
% 3. 특히, pre-training scale 이 커질수록 더 성능 향상폭이 증가함 -- naive fine-tuning 과는 달리
% 4. 마지막으로, SWAD 랑 결합해서 더 좋아질 수 있셤
% 주장하는 바가 여러가지라서 figure 로 통합한다고 해서 분량이 줄 거 같지는 않음...
In our experiment, we observe that the naive fine-tuning of a larger pre-trained model can fail to provide better performances, even though the larger pre-trained model is trained with more data and domains. For example, ERM with the ResNet pre-trained on ImageNet (trained with 1.3M images) shows 64.2\% of averaged accuracy, while ERM with the ViT pre-trained on CLIP (trained with 400M image-caption pairs) shows 61.1\%.
On the other hand, we show that our method can significantly improve the average DG performances with backbone models at different scales, \eg, ImageNet pre-trained ResNet (64.2\% $\rightarrow$ 65.9\%), 400M image-text pre-trained ViT (CLIP) \cite{radford2021clip} (61.1\% $\rightarrow$ 73.7\%) and Instagram 3.6B pre-trained RegNet (SWAG) \cite{singh2022swag} (68.0\% $\rightarrow$ 74.1\%). Especially, we observe that the pre-trained knowledge by larger pre-trained models, such as SWAG and CLIP, is more effective to learn domain generalized features than the ImageNet pre-trained model: \method{} with the ViT pre-trained on CLIP outperforms \method{} with the ResNet pre-trained on ImageNet in contrast to the naive fine-tuning.
Furthermore, our feature-level regularization method is easily combined with the existing parameter space ensemble methods \cite{cha2021swad, wortsman2021robust} (74.1\% $\rightarrow$ \textbf{77.3\%} average DG accuracy by combining with SWAD \cite{cha2021swad} and pre-trained RegNet).

Our contribution is as follows: (1) We re-formulate the DG objective by mutual information with the oracle model. Then, we approximate the oracle by a large pre-trained model to derive a tractable approximation of the target objective. We propose \fullname{} (\method{}) to solve our objective. (2) We analyze the pre-trained models in terms of the MI with the oracle model. Our analysis shows that naive fine-tuning of pre-trained models can harm the MI with the oracle, on the other hand, \method{} shows high MI with the oracle.
% Our analysis provides a weak link between the bad out-of-domain generalizability by naive fine-tuning and mutual information with oracle model. 
(3) We compare \method{} with state-of-the-art DG methods on DomainBed. \method{} outperforms all methods in all settings, including varying optimizers and pre-trained models. We also provide extensive analysis to understand \method{}. For example, we observe that \method{} shows stronger DG performances with larger pre-trained models, such as SWAG \cite{singh2022swag} or CLIP \cite{radford2021clip}.


% We begin with a re-formulation of domain generalization in the lens of mutual information between the oracle model and our target model, and deriving a tractable variational bound by approximating the oracle model by a large pre-trained model. We empirically validate that our approximation of the oracle model by showing the actual mutual information quantity between the oracle model and the pre-trained models. Finally, we present our results and discussions related to the choice of pre-trained models.